\section*{अध्याय \ref{ch:g11:lithosphere-asthenosphere} --- स्थलमंडल और दुर्बलतामंडल}

\begin{solution}{exo:g11:lithosphere-asthenosphere:1}
\omterm{def:g11:lithosphere-asthenosphere:lithosphere}{स्थलमंडल} पृथ्वी का कठोर बाहरी कवच है, जो भूपर्पटी और
सबसे ऊपरी प्रावार से बना है; \omterm{def:g11:lithosphere-asthenosphere:lithosphere}{दुर्बलतामंडल} उसके नीचे का अधिक गरम, अधिक दुर्बल प्रावार है,
जो लंबे समय में धीरे-धीरे बहता है। महासागरीय भूपर्पटी, महाद्वीपीय
भूपर्पटी और सबसे ऊपरी प्रावार \omterm{def:g11:lithosphere-asthenosphere:lithosphere}{स्थलमंडल} के अंग हैं; बाह्य क्रोड
नहीं।
\end{solution}

\begin{solution}{exo:g11:lithosphere-asthenosphere:2}
हिमशैल तब तक धँसता है जब तक वह अपने भार के बराबर जल न हटा दे:
जल के नीचे का भाग $910/1025 \approx 0.888$ है, लगभग \qty{89}{\percent}।
ऊपर का भाग लगभग \qty{11.2}{\percent} है, इसलिए कुल ऊँचाई
$20/0.112 \approx \qty{180}{m}$ है, जिसमें से लगभग \qty{160}{m} जल के
नीचे है।
\end{solution}

\begin{solution}{exo:g11:lithosphere-asthenosphere:3}
$15 + 26 \times 3 = \qty{93}{\celsius}$। \qty{100}{km} पर वही प्रवणता
\qty{2600}{\celsius} से अधिक देती, जो प्रावार की चट्टान के गलनांक से बहुत
ऊपर है; पर \omterm{def:g11:seismic-structure:body-waves}{S तरंगें} ऊपरी प्रावार को पार करती हैं, जो इसलिए ठोस है।
गहराई के साथ प्रवणता को बहुत छोटा होना चाहिए।
\end{solution}

\begin{solution}{exo:g11:lithosphere-asthenosphere:4}
समुद्र-तल से 0 से \qty{500}{m} ऊपर की पट्टी, सतह के लगभग
\qty{16}{\percent} के साथ, और समुद्र-तल से 4000 से \qty{4500}{m}
नीचे की पट्टी, लगभग \qty{11}{\percent} के साथ (4500 से
\qty{5000}{m} नीचे की पट्टी लगभग उतनी ही बड़ी है)।
\end{solution}

\begin{solution}{exo:g11:lithosphere-asthenosphere:5}
$P = 0$ पर: \qty{1121}{\celsius}। \qty{3}{GPa} पर: $-5.104 \times 9 +
132.899 \times 3 + 1120.661 \approx -45.9 + 398.7 + 1120.7 \approx
\qty{1473}{\celsius}$।
\end{solution}

\begin{solution}{exo:g11:lithosphere-asthenosphere:6}
\qty{50}{km} पर लगभग \qty{730}{\celsius} (\omterm{def:g11:lithosphere-asthenosphere:lithosphere}{स्थलमंडल}),
\qty{150}{km} पर लगभग \qty{1370}{\celsius} और \qty{250}{km} पर
\qty{1410}{\celsius} (\omterm{def:g11:lithosphere-asthenosphere:lithosphere}{दुर्बलतामंडल})। तापमान 50 और
\qty{150}{km} के बीच सबसे अधिक बढ़ता है, \qty{600}{\celsius} से अधिक, क्योंकि उस
अंतराल का अधिकांश भाग \omterm{def:g11:lithosphere-asthenosphere:lithosphere}{स्थलमंडल} में है; 150 और \qty{250}{km} के बीच वह
केवल लगभग \qty{45}{\celsius} बढ़ता है।
\end{solution}

\begin{solution}{exo:g11:lithosphere-asthenosphere:7}
मूल: $3 \times 2830/550 \approx \qty{15.4}{km}$। अतिरिक्त मोटाई:
$3 \times 3380/550 \approx \qty{18.4}{km}$, यानी श्रेणी के \qty{3}{km}
और उसका मूल।
\end{solution}

\begin{solution}{exo:g11:lithosphere-asthenosphere:8}
भूपर्पटी $55 - 41.1 = \qty{13.9}{km}$ अधिक मोटी है, और $\Delta H = h
\times 3380/550$, इसलिए $h = 13.9 \times 550/3380 \approx \qty{2.3}{km}$।
\end{solution}

\begin{solution}{exo:g11:lithosphere-asthenosphere:9}
कमी: $(4.470 - 4.419)/4.470 \approx \qty{1.1}{\percent}$। उछाल:
$(4.644 - 4.419)/4.419 \approx \qty{5.1}{\percent}$। गहराई के साथ चट्टान
अधिकाधिक दबती है, जिससे भूकंपीय तरंगें तेज़ होती हैं; कमी का अर्थ है
कि कोई और बात भारी पड़ती है: चट्टान गरम है और पिघलने के निकट है, और अधिक दुर्बल है।
\end{solution}

\begin{solution}{exo:g11:lithosphere-asthenosphere:10}
(1) $(1345 - 610)/60 \approx \qty{12}{\celsius}$ प्रति किलोमीटर। (2) यह
सतह की प्रवणता का लगभग आधा है। ऊष्मा का एक भाग ऊपरी
भूपर्पटी के भीतर बनता है; बाहर निकलने के लिए उस ऊष्मा को भी ऊपरी परतें पार करनी पड़ती हैं, जिससे
सतह के पास प्रवणता अधिक तीव्र हो जाती है। भूपर्पटी के नीचे केवल अधिक गहराई से आने वाली
ऊष्मा चट्टान को पार करती है, और प्रवणता कम होती है।
\end{solution}

\begin{solution}{exo:g11:lithosphere-asthenosphere:11}
\qty{200}{km} पर: $1300 + 90 = \qty{1390}{\celsius}$; \qty{400}{km} पर:
$1300 + 180 = \qty{1480}{\celsius}$। वृद्धि \qty{200}{km} में \qty{90}{\celsius}
है। \omterm{def:g11:lithosphere-asthenosphere:lithosphere}{स्थलमंडल} के आर-पार तापमान केवल \qty{100}{km} में
$1345 - 15 = \qty{1330}{\celsius}$ बढ़ता है: प्रति किलोमीटर लगभग तीस
गुना अधिक।
\end{solution}

\begin{solution}{exo:g11:lithosphere-asthenosphere:12}
(1) $118 + 12.3\,z = 1300 + 0.45\,z$ से $z = 1182/11.85 \approx
\qty{100}{km}$ मिलता है। (2) $118 + 21.5\,z = 1300 + 0.45\,z$ से $z = 1182/21.05
\approx \qty{56}{km}$ मिलता है। (3) अधिक गरम पश्चिम का \omterm{def:g11:lithosphere-asthenosphere:lithosphere}{स्थलमंडल} बहुत पतला है,
जैसा मापा गया है (सबसे पतले स्थान पर लगभग \qty{40}{km}); महाद्वीप के मध्य के बहुत मोटे
\omterm{def:g11:lithosphere-asthenosphere:lithosphere}{स्थलमंडल} के लिए और भी छोटी प्रवणता चाहिए।
\end{solution}

\begin{solution}{exo:g11:lithosphere-asthenosphere:13}
(1) प्रस्ताव के अनुसार $\Delta H = h\,\rho_m/(\rho_m - \rho_c)$, इसलिए $h =
\Delta H\,(\rho_m - \rho_c)/\rho_m$। (2) $550/3380 \approx 0.16$: भूपर्पटी की किसी भी अतिरिक्त
मोटाई का लगभग छठा भाग ऊँचाई के रूप में दिखता है, पाँच-छठा भाग मूल के रूप में। (3)
हिमशैल के लिए जल से ऊपर का भाग $(1025 - 910)/1025 \approx 0.11$ है:
वही नियम, (तरल का घनत्व $-$ तैरती वस्तु का घनत्व)/तरल का
घनत्व।
\end{solution}

\begin{solution}{exo:g11:lithosphere-asthenosphere:14}
(1) $1489 - 1345 = \qty{144}{\celsius}$। (2) थोड़ा पिघला पदार्थ चट्टान के कणों के बीच
पतली परतों में इकट्ठा होता है। तब कण आपस में मज़बूती से
जकड़े नहीं रहते: चट्टान अपरूपण का कम प्रतिरोध करती है, इसलिए \omterm{def:g11:seismic-structure:body-waves}{S तरंगें} धीमी होती हैं, और वह
कहीं अधिक आसानी से विकृत हो सकती है और बह सकती है।
\end{solution}

\begin{solution}{exo:g11:lithosphere-asthenosphere:15}
(1) नहीं। स्थल को समतल करने से महाद्वीपीय भूपर्पटी थोड़ी पतली हो जाती,
पर वह फिर भी मोटी और हल्की रहती, और महासागरीय तली की पतली, घनी भूपर्पटी से
अब भी बहुत ऊँचे तैरती; महासागर अब भी
महाद्वीपों के बीच की द्रोणियाँ भरते। (2) भूपर्पटी दो प्रकार की है, अलग-अलग
मोटाई और घनत्व की: एक मोटी, हल्की महाद्वीपीय भूपर्पटी और एक पतली, घनी
महासागरीय भूपर्पटी।
\end{solution}

\begin{solution}{pb:g11:lithosphere-asthenosphere:1}
\textbf{1.} $5.01 - 0.67 = \qty{4.34}{km}$।

\textbf{2.} $(3380 - 2830)/3380 \approx \qty{16}{\percent}$।

\textbf{3.} भिन्न $2830/3380 \approx 0.84$ नीचे होता;
हिमशैल थोड़ा अधिक गहरा धँसता है, $910/1025 \approx 0.89$, क्योंकि बर्फ़ और जल के घनत्व
भूपर्पटी और प्रावार के घनत्वों की अपेक्षा एक-दूसरे के अधिक निकट हैं।

\textbf{4.} भूकंपीय तरंग के सेकंडों में, या किसी भूकंप की अवधि में,
प्रावार ठोस है; केवल बहुत लंबे समय में ही
\omterm{def:g11:lithosphere-asthenosphere:lithosphere}{दुर्बलतामंडल} बहता है और भूपर्पटी को तैरने देता है।

\textbf{5.} सबसे गहरे \omterm{def:g11:seismic-structure:moho}{मोहो} के नीचे दोनों स्तंभ उन्हीं गहराइयों पर एक ही प्रावार-चट्टान
से बने हैं, और वहाँ उनका भार समान है; सारे अंतर
उस गहराई के ऊपर हैं और उन्हें सम्मिलित करना ही होगा।

\textbf{6.} औसत महाद्वीप: समुद्र-तल से \qty{0.67}{km} ऊपर से
\qty{40.4}{km} नीचे तक भूपर्पटी, फिर \qty{70}{km} तक प्रावार।
पठार: समुद्र-तल से \qty{5.01}{km} ऊपर से नीचे मूल जितने अधिक गहरे
\omterm{def:g11:seismic-structure:moho}{मोहो} तक भूपर्पटी, फिर \qty{70}{km} तक प्रावार।

\textbf{7.} पठार के स्तंभ में $h + r$ अधिक भूपर्पटी है, महाद्वीप के स्तंभ में
उसके बजाय मोटाई $r$ तक प्रावार: $\rho_c\,(h + r) = \rho_m\,r$,
इसलिए $r = h\,\rho_c/(\rho_m - \rho_c)$ और $\Delta H = h + r =
h\,\rho_m/(\rho_m - \rho_c)$।

\textbf{8.} $3380/550 \approx 6.15$।

\textbf{9.} $\Delta H = 4.34 \times 6.15 \approx \qty{26.7}{km}$; मूल
$26.7 - 4.34 \approx \qty{22.3}{km}$ है।

\textbf{10.} औसत महाद्वीप का \omterm{def:g11:seismic-structure:moho}{मोहो} समुद्र-तल से $41.1 - 0.67 \approx
\qty{40.4}{km}$ नीचे है; पठार का $40.4 + 22.3 \approx
\qty{62.8}{km}$ नीचे।

\textbf{11.} $41.1 + 26.7 \approx \qty{67.8}{km}$।

\textbf{12.} एक वर्ग मीटर के ऊपर एक किलोमीटर चट्टान का द्रव्यमान
$\rho \times \qty{1000}{m^3}$ है: भूपर्पटी का \num{2830} टन या प्रावार का \num{3380}
टन। महाद्वीप: $41.1 \times 2830 + (70 - 40.43) \times 3380
\approx \num{116300} + \num{99900} \approx \num{216300}$ टन। पठार:
$67.77 \times 2830 + (70 - 62.76) \times 3380 \approx \num{191800} +
\num{24500} \approx \num{216300}$ टन। दोनों बराबर हैं।

\textbf{13.} $41.1 - 0.67 \times 6.15 \approx \qty{37.0}{km}$।

\textbf{14.} $67.8/41.1 \approx 1.65$ गुना।

\textbf{15.} $3300/470 \approx 7.02$, $\Delta H \approx 4.34 \times 7.02
\approx \qty{30.5}{km}$, और लगभग \qty{71.6}{km} की भूपर्पटी। प्रावार के घनत्व में
लगभग \qty{2}{\percent} का परिवर्तन मोटाई को
लगभग \qty{4}{km} बदल देता है: परिणाम घनत्वों पर स्पष्ट रूप से निर्भर करता है, पर
\qty{70}{km} के निकट रहता है।

\textbf{16.} मॉडल लगभग \qty{68}{km} देता है, भूकंपीय तरंगें लगभग
\qty{70}{km}: वे कुछ प्रतिशत के भीतर सहमत हैं।

\textbf{17.} $h = 3.23 - 0.67 = \qty{2.56}{km}$, $\Delta H = 2.56 \times
6.15 \approx \qty{15.7}{km}$, और लगभग \qty{56.8}{km} की भूपर्पटी, जो मापी गई मोटाई से
लगभग \qty{7}{km} अधिक है।

\textbf{18.} पतली। जिस स्तंभ का ऊपरी भाग औसत से हल्का है, उसका भार
उतनी ही मोटाई के लिए कम होता है, इसलिए वह अधिक ऊँचे तैरता है: हल्की भूपर्पटी को
दी हुई ऊँचाई पर खड़े रहने के लिए छोटा मूल चाहिए। इसलिए वास्तविक भूपर्पटी
मॉडल की भूपर्पटी से पतली हो सकती है, जैसा मापा गया है।

\textbf{19.} लंबे समय में \omterm{def:g11:lithosphere-asthenosphere:lithosphere}{दुर्बलतामंडल} बहता है: वह भारी स्तंभ के नीचे से
हटता है और हल्के स्तंभ के नीचे लौट आता है, जब तक स्तंभ संतुलित न हो जाएँ।

\textbf{20.} \textbf{तिब्बत के नीचे की भूपर्पटी लगभग \qty{70}{km} मोटी है}, जो
औसत महाद्वीप के \qty{41}{km} से बहुत अधिक है: पठार पाँच
किलोमीटर ऊँचा इसलिए है कि उसकी मोटी, हल्की भूपर्पटी अधिक घने प्रावार पर तैरती है,
\qty{20}{km} से अधिक के मूल के साथ।
\end{solution}
